\subsection{Regimes}\label{sec:s2-regime}
Table~\ref{tab:s2regime} and Fig.~\ref{fig:s2regime} break the utility down by disturbance level, time-window factor and
latency factor. Per configuration and without correction for multiple testing, \mname{} without fallback was better than
ALNS with 40\,ms in \val{KWinForty} of the \val{KNConfigs} configurations and worse in \val{KLossForty}, better than the
constant buffer in \val{KWinBuf} and worse in \val{KLossBuf}, and better than the conventional planner in \val{KWinConv}
and worse in \val{KLossConv}. The advantage over ALNS with 40\,ms was \val{KDFortyKTen}\,pp, \val{KDFortyKFifty}\,pp and
\val{KDFortyKTwoHundred}\,pp for $\kappa_B=10$, 50 and 200, and the advantage over online prediction stayed between
\val{KDOnlKTen}\,pp and \val{KDOnlKTwoHundred}\,pp.

The advantage was largest with the original windows. Under high disturbances, \mname{} without fallback served
\val{KDFortyWOneHigh}\,pp more reward than ALNS with 40\,ms and \val{KDBufWOneHigh}\,pp more than the constant buffer.
The losses all occurred under high disturbances with halved windows. In this regime, \mname{} without fallback reached
\val{KUPropWHalfHigh}\%, \val{KDBufWHalfHigh}\,pp relative to the constant buffer and \val{KDFortyWHalfHigh}\,pp relative
to ALNS with 40\,ms. It served \val{KServedPropWHalfHigh} tasks per mission, of which \val{KLatePropWHalfHigh} late,
against \val{KServedBufWHalfHigh} and \val{KLateBufWHalfHigh} for the constant buffer. The chance-constrained model was
thus too cautious when the windows were short compared with the predicted finish-time error. The effective deadline of
many tasks then probably falls before the earliest possible arrival, so the planner gives up tasks that a less cautious
plan still serves in time. The value of $z$ was selected per disturbance level on development missions, averaged over
both window widths. Selecting it also per window width, or adapting it online, is a natural extension.

\inputtable{tabm_s2_regime}

\begin{figure}[t]
\centering
\fig{fig_s2_regime}{\linewidth}
\caption{Study 2: utility per regime against the latency factor, mean over \val{KNMissions} missions per point.}\label{fig:s2regime}
\end{figure}

\subsection{Tier C: PX4 software-in-the-loop with Gazebo}\label{sec:Ss2-tierc}
All \val{KCRuns} held-out flights of study~2 were completed, that is, \val{KCN} missions with four methods each. One run
failed to start and was repeated. As specified in the protocol, \val{KCInvalid} runs with a real-time factor outside
1.6 to 2.2 were repeated once with the same seed, and the original runs were kept. The laptop was also used for other
work during the campaign, and many repeats ran slower than the original runs. Of the \val{KCRuns} analysed runs,
\val{KCValid} lie in the validity range. The real-time factor was \val{KCRTF} on average, with quartiles of
\val{KCRTFq} and a minimum of \val{KCRTFmin}. Two sensitivity analyses that were not planned in the protocol therefore repeat the
comparisons with the original runs instead of the repeats, and with only those pairs of runs that both lie in the
validity range.

Table~\ref{tab:s2tierc} summarises the results. \mname{} without fallback reached a utility of \val{KCUProp}\%, against
\val{KCUForty}\% for ALNS with 40\,ms, \val{KCUBuf}\% for the constant buffer and \val{KCUConv}\% for the conventional
planner. The paired differences were \val{KCDForty}\,pp to ALNS with 40\,ms and \val{KCDConv}\,pp to the conventional
planner, with Holm-adjusted $p$-values of \val{KCPForty} and \val{KCPConv}. The fourth hypothesis is therefore not
supported. The secondary comparison with the constant buffer gave \val{KCDBuf}\,pp. The sensitivity analyses lead to
the same conclusion. With the original runs, the differences were \val{KCDFortyOrig}\,pp and \val{KCDConvOrig}\,pp.
With valid runs only, they were \val{KCDFortyValid}\,pp on \val{KCNFortyValid} missions and \val{KCDConvValid}\,pp on
\val{KCNConvValid} missions. No interval in the three analyses excludes zero.

\inputtable{tabm_s2_tierc}

The mechanism seen in tier B reappears in the counts of late and skipped tasks. \mname{} without fallback served
\val{KCLateProp} tasks late per mission, against \val{KCLateForty} for ALNS with 40\,ms, which is \val{KCLateRedForty}\%
fewer. It skipped no task at departure, whereas ALNS with 40\,ms skipped \val{KCSkipForty} and the conventional planner
\val{KCSkipConv} tasks per mission. It computed \val{KCMsProp}\,ms per call, against \val{KCMsForty}\,ms for the fixed
budget. The conventional planner called the search \val{KCCallsConv} times per mission instead of about
\val{KCCallsProp}, as in tier B.


The same \val{KCN} missions were also run in tier B with the rescaled parameters of tier C. As in study~1, tier B
predicted the utility of a flight closely, with a correlation of $r=$ \val{KBCr} over \val{KBCn} runs. On these
missions, however, even tier B did not resolve the differences between the methods (Fig.~\ref{fig:s2bc}). Its estimates
were \val{KBgzDForty}\,pp to ALNS with 40\,ms, \val{KBgzDBuf}\,pp to the constant buffer and \val{KBgzDConv}\,pp to the
conventional planner. The tier C intervals overlap them. The per-mission difference to ALNS with 40\,ms had a standard
deviation of \val{KCSdForty}\,pp in tier C. Confirming the tier B estimate of \val{KCEffForty}\,pp on these missions with
a power of 80\% would require about \val{KCNPowerForty} missions per method. Tier C thus shows that the proposed
replanner runs on an autopilot and a physics engine with the same computation and fewer late tasks, but it cannot
confirm the size of the gain measured in tier B.

\begin{figure}[t]
\centering
\fig{fig_s2_bc}{0.62\linewidth}
\caption{Study 2, tier B against tier C on the same \val{KCN} missions: \mname{} without fallback minus each baseline, with 95\% intervals.}
\label{fig:s2bc}
\end{figure}
